% ===================================================================
%  CADRO N.º 11 — A cidade e os seus monumentos. — O incendio.
%  Traduction 2026 d'apres texto/io, controlee sur texto/fr, texto/es
%  et texto/pt. Consigne : voir texto/gl/10-cadro-01.tex.
%
%  OUVERTURE DE LA TROISIEME SERIE : « TERCEIRA SERIE ».
%
%  LE TABLEAU DE LA VILLE EST LE PLUS INTERNATIONAL DE LA COLONNE,
%  plus encore que celui de la mer, et pour une raison qui acheve la
%  demonstration commencee au tableau 2 : la ville du XIXe siecle a
%  ete BATIE d'apres un plan, et le plan avait un vocabulaire.
%  « Bulevar », « prefectura », « liceo », « museo », « biblioteca »,
%  « tranvía », « hospital », « telégrafo » sont les memes mots dans
%  les trois langues, et dans les vingt-cinq autres colonnes du
%  volume.
%
%  LE BLOC c2-02-1 A COUTE UNE INVERSION, ET LES DEUX VOISINES L'ONT
%  PAYEE AUSSI. L'ido ecrit « plu bone kam la soldati (57) o la
%  jendarmi (63), retroirigas la kuriozi (56) » : le terme de
%  comparaison AVANT l'objet. Aucune des trois langues romanes ne
%  range ainsi — elles rejettent la comparaison apres le verbe et
%  l'objet apres elle — et l'espagnol comme le roumain ont du refaire
%  la phrase de la meme facon : faire de la comparaison le predicat
%  (« lo hacen mejor que... para hacer retroceder », « se descurcă mai
%  bine decât... împingând înapoi »). Le galicien fait comme eux :
%  « fano mellor ca os soldados (57) ou os gardas (63) para faceren
%  recuar os curiosos (56) ». Le meme bloc avait coute une inversion
%  a l'irlandais, mais pour une autre raison et avec un autre remede
%  — voir le message de commit de son tableau 11.
%
%  ET C'EST LE CONTRAIRE EXACT DU TABLEAU 7. La ferme donnait un mot
%  galicien sur deux ; la ville n'en donne presque aucun. Entre les
%  deux, la meme regle : la ferme s'apprend dehors et sans maitre, la
%  ville s'apprend a l'ecole, dans l'administration et dans le
%  journal. Ce n'est pas le sujet qui varie, c'est l'ECOLE — et la
%  ville EST une ecole.
% ===================================================================

%%K t11-apar-1 apar
\VUsaut{2.0mm}
\VUcentre{13.2pt}{90}{TERCEIRA SERIE}
\VUsaut{3.0mm}
\VUfilet{16mm}
\VUsaut{5.6mm}
\VUtitre{15.0pt}{60}{CADRO N.º 11}
\VUsaut{1.4mm}
\VUpk{11.4pt}{40}{A cidade e os seus monumentos. --- O incendio.}
\VUsaut{2.2mm}

%%K t11-tit-1 sub
\VUpk{10.2pt}{40}{O arrabalde.}
\VUsaut{3.0mm}

%%K t11-c1-01-1 p
1. --- Vivo nunha cidade grande situada a cen quilómetros do Océano e que, aínda así,
é un \VUgras{porto} \textsuperscript{(1)} de primeira orde. O río que a bordea ten
catrocentos metros de largo e forma unha rada moi fermosa.

%%K t11-c1-02-1 p
2. --- Toda a parte da cidade que está nos \VUgras{cais} \textsuperscript{(2)} ten un
aspecto totalmente mariñeiro e industrial, e forma un \VUgras{arrabalde}
\textsuperscript{(3)} habitado só por mariñeiros e obreiros. Estes traballan nas
\VUgras{fábricas} \textsuperscript{(4)} e na \VUgras{fábrica de gas}
\textsuperscript{(5)}, cuxa alta \VUgras{cheminea} \textsuperscript{(6)} e cuxo
\VUgras{gasómetro} se ven de moi lonxe.

%%K t11-c1-03-1 p
3. --- Na Idade Media a cidade estaba fortificada, e aínda se atopan algúns restos das
súas murallas, en particular a \VUgras{porta} \textsuperscript{(8)} do vello
\VUgras{castelo fortificado} \textsuperscript{(7)}, flanqueado polas súas dúas
\VUgras{torres} \textsuperscript{(9)}, que hoxe serven de \VUgras{cárcere.}

%%K t11-c1-04-1 p
4. --- En todo ese \VUgras{barrio,} que ten un aspecto moi vello, un só edificio é
relativamente moderno: é a \VUgras{aduana} \textsuperscript{(10)}. Está situada entre
o cais e a \VUgras{rúa pequena} \textsuperscript{(11)} que leva á vella \VUgras{casa
do concello} \textsuperscript{(12)}. Recoñécese doadamente ese último monumento polo
\VUgras{campanario} \textsuperscript{(13)} xigantesco que domina a cidade antiga.

%%K t11-c1-05-1 p
5. --- Do outro lado da rúa \VUgras{ergue} un edificio construído no mesmo estilo ca a
aduana: é a \VUgras{Bolsa} \textsuperscript{(14)}. Unha das súas fachadas dá a unha
pequena praza na que está a \VUgras{prefectura} \textsuperscript{(15)}. Efectivamente,
a nosa cidade, sen ser unha capital, é aínda así unha importante cabeceira de
departamento.

%%K t11-tit-2 sub
\VUsaut{5.0mm}
\VUpk{10.2pt}{40}{A parte máis alta da cidade.}
\VUsaut{3.4mm}

%%K t11-c1-06-1 p
6. --- A guarnición, bastante numerosa, está aloxada en amplos \VUgras{cuarteis}
\textsuperscript{(16)} moi salubres; esténdense ata a reixa do \VUgras{parque
municipal} \textsuperscript{(17)}. O \VUgras{bulevar} \textsuperscript{(19)} que leva
a ese parque pasa detrás da \VUgras{caixa de aforros} \textsuperscript{(18)} e vai
dar á praza da \VUgras{catedral} \textsuperscript{(20)}.

%%K t11-c1-07-1 p
7. --- Todos eses monumentos espállanse en anfiteatro nun outeiro pequeno onde está
tamén o noso amplo \VUgras{instituto} \textsuperscript{(21)}.

%%K t11-c1-07-2 p
Se un baixa por unha vía algo en costa, que se xunta á Rúa Grande, chégase ao
\VUgras{Palacio de Xustiza} \textsuperscript{(22)}, diante do cal se estende un
agradable \VUgras{xardín público} \textsuperscript{(26)}. Esa \VUgras{praza
axardinada} non é moi grande, pero pon no medio da cidade un oasis de verdura e
frescura. Por iso é moi frecuentada polos habitantes da cidade, que gustan de vir
sentar baixo o toldo do \VUgras{café} \textsuperscript{(25)} para escoitar os músicos
soldados que os xoves e os domingos tocan no \VUgras{quiosco}
\textsuperscript{(27)} posto entre os céspedes e os macizos de flores.

%%K t11-c1-08-1 p
8. --- Nun deses céspedes érguese unha \VUgras{estatua} \textsuperscript{(28)} de
bronce, monumento erguido á memoria dun dos nosos conveciños que fixo grandes servizos
á súa cidade natal.

%%K t11-tit-3 sub
\VUsaut{4.0mm}
\VUpk{10.2pt}{40}{Os medios de transporte.}
\VUsaut{3.4mm}

%%K t11-c1-09-1 p
9. --- Ata agora a nosa cidade non está aínda alumada coa luz eléctrica; só ten
\VUgras{lámpadas de gas} \textsuperscript{(29)}. Pero \VUgras{a prensa local} fai
campaña para que ese sistema de alumado sexa mellorado, \VUgras{como} xa se
perfeccionou \VUgras{o sistema de tracción dos} tranvías. \VUgras{Os} vellos coches,
\VUgras{tirados} por \VUgras{cabalos,} foron practicamente substituídos por
\VUgras{tranvías eléctricos} \textsuperscript{(31)} que transportan axiña os viaxeiros
dun extremo da \VUgras{cidade ao outro,} por un prezo moi \VUgras{barato.}

%%K t11-c1-10-1 p
10. --- Preto do Palacio de Xustiza está o \VUgras{Banco} \textsuperscript{(32)} no
que se descontan os valores comerciais. Os \VUgras{ordenanzas do banco}
\textsuperscript{(33)} van cobrar a domicilio as débedas.

%%K t11-tit-4 sub
\VUsaut{4.0mm}
\VUpk{10.2pt}{40}{Arte e instrución.}
\VUsaut{3.4mm}

%%K t11-c1-11-1 p
11. --- O belo edificio que se ergue preto de alí é o \VUgras{museo.} Contén á vez a
\VUgras{biblioteca} \textsuperscript{(34)} da cidade, as belas \VUgras{coleccións de
ciencias naturais} \textsuperscript{(35)} (Museo de Historia Natural) e unha graciosa
\VUgras{galería de pintura} \textsuperscript{(36)} que constitúe unha espléndida
\VUgras{exposición} permanente.

%%K t11-c1-11-2 p
Ábrese ao público dúas veces por semana, pero os forasteiros poden visitalo calquera
día dando unha propina ao \VUgras{garda} \textsuperscript{(37)}. Os \VUgras{pintores}
\textsuperscript{(38)} veñen alí traballar. Vense \VUgras{diante} do seu cabalete, coa
\VUgras{paleta} na man, copiando os cadros famosos.

%%K t11-c1-12-1 p
12. --- No alto, detrás do museo, albíscase a \VUgras{cúpula} \textsuperscript{(40)}
do \VUgras{hospital} \textsuperscript{(39)} e os edificios da \VUgras{escola normal}
\textsuperscript{(41)} na que se formaban os mestres das nosas escolas municipais. Na
praza do Teatro hai unha \VUgras{oficina de correos} \textsuperscript{(43)} instalada
nunha \VUgras{casa particular} \textsuperscript{(42)}.

%%K t11-tit-5 sub
\VUsaut{5.0mm}
\VUpk{11.4pt}{40}{O incendio.}
\VUsaut{3.0mm}

%%K t11-tit-6 sub
\VUpk{10.2pt}{40}{Berro de incendio.}
\VUsaut{3.4mm}

%%K t11-c2-01-1 p
1. --- Un rumor arrepiante espállase pola cidade: acaba de estalar un incendio no
teatro! No intre en que chego á \VUgras{praza} \textsuperscript{(96)}, a multitude
está xa amoreada no \VUgras{beirarrúa} \textsuperscript{(46)} e na \VUgras{calzada}
\textsuperscript{(47)}. Os \VUgras{empregados de correos} \textsuperscript{(44)} e os
\VUgras{carteiros} \textsuperscript{(45)} saen da oficina. Un \VUgras{condutor de
tranvía} \textsuperscript{(48)} tivo que parar o seu coche para deixar pasar un
\VUgras{coche de ambulancias} municipais \textsuperscript{(50)}, que chega cun
\VUgras{cirurxián} \textsuperscript{(52)} e un \VUgras{enfermeiro}
\textsuperscript{(51)} para atender un \VUgras{ferido} \textsuperscript{(53)} traído
nunha \VUgras{padiola} \textsuperscript{(54)} preto do \VUgras{quiosco de xornais}
\textsuperscript{(55)}.

%%K t11-c2-02-1 p
2. --- Unha compañía de infantaría, que volvía do exercicio, parou para asegurar o
servizo de orde cos \VUgras{gardas municipais a cabalo} \textsuperscript{(61)}. Os
\VUgras{cabaleiros,} cuxos cabalos se encabritan, fano mellor ca os \VUgras{soldados}
\textsuperscript{(57)} ou os \VUgras{gardas} \textsuperscript{(63)} para faceren
recuar os \VUgras{curiosos} \textsuperscript{(56)}. Ademais os gardas están moi ocupados. Un
deles
indícalles aos \VUgras{axentes de policía} \textsuperscript{(64)} onde cómpre levar
outra \VUgras{vítima} \textsuperscript{(65)}, un valente bombeiro caído dunha
escaleira.

%%K t11-c2-03-1 p
3. --- \VUgras{O oficial} \textsuperscript{(66)} precisa o lugar onde se debe colocar
a \VUgras{bomba de vapor} \textsuperscript{(67)} que chega. Mentres agarda por esa
potente máquina, mandou pór en batería unha \VUgras{bomba de man}
\textsuperscript{(68)}, cuxo longo \VUgras{tubo} \textsuperscript{(69)} lanza auga a
torrentes sobre o lume. Seis bombeiros manóbrana, mentres o seu compañeiro, que ten a
\VUgras{lanza} \textsuperscript{(70)}, dirixe o \VUgras{chorro}
\textsuperscript{(71)} ao centro do incendio.

%%K t11-c2-04-1 p
4. --- Do outro lado do teatro apoiouse a grande \VUgras{escaleira}
\textsuperscript{(72)}. Permítelles aos bombeiros achegárense ás lapas que xorden entre
remuíños de \VUgras{fume} negro \textsuperscript{(75)}.

%%K t11-tit-7 sub
\VUsaut{5.0mm}
\VUpk{10.2pt}{40}{Fazañas valentes.}
\VUsaut{3.4mm}

%%K t11-c2-05-1 p
5. --- \VUgras{O incendio} \textsuperscript{(74)} naceu no vestíbulo do
\VUgras{teatro} \textsuperscript{(73)}, mentres se ensaiaba \VUgras{a ópera} cuxa
primeira \VUgras{representación} había ser mañá. Por sorte só había poucos
\VUgras{espectadores} \textsuperscript{(76)} na sala. Os coitados refuxiáronse no
\VUgras{balcón} \textsuperscript{(77)}, onde uns valentes \VUgras{bombeiros}
\textsuperscript{(86)} veñen socorrelos con escaleira e cordas.

%%K t11-c2-06-1 p
6. --- Baixo o \VUgras{peristilo} \textsuperscript{(78)}, onde o \VUgras{cartel}
\textsuperscript{(79)} anuncia a representación, vense fuxir os \VUgras{músicos}
\textsuperscript{(81)} da \VUgras{orquestra} \textsuperscript{(80)}, levando os seus
instrumentos: \VUgras{violín} \textsuperscript{(81)}, \VUgras{frauta}
\textsuperscript{(82)}, \VUgras{violonchelo} \textsuperscript{(83)},
\VUgras{trombón} \textsuperscript{(84)}, \VUgras{tambor} \textsuperscript{(85)}, etc.
Tenta salvarse unha parte do \VUgras{mobiliario} \textsuperscript{(87)}.

%%K t11-c2-07-1 p
7. --- \VUgras{Os actores} \textsuperscript{(89)} e \VUgras{as actrices}
\textsuperscript{(88)}, \VUgras{cantantes} e \VUgras{cantatrices,} tiveron que fuxir
co seu \VUgras{traxe} \textsuperscript{(90)} de teatro. O tenor explícalle a un
\VUgras{xornalista} \textsuperscript{(92)} como sucedeu. \VUgras{O reporteiro}
acudiu axiña desde o primeiro intre coas autoridades: o \VUgras{prefecto}
\textsuperscript{(91)}, o \VUgras{alcalde} \textsuperscript{(93)}, o
\VUgras{comisario de policía} \textsuperscript{(94)} e un \VUgras{maxistrado}
\textsuperscript{(95)}, que abrirán unha investigación para xulgar as
responsabilidades.
